\begin{proposition}[7.9.3]
\label{IV.2.s7.9.3-fr}
Sea $A$ un anillo local noetheriano reducido, y supongamos que se puede
resolver $\operatorname{Spec}(A)$. Entonces las fibras del morfismo
canónico $\operatorname{Spec}(\widehat A)\to\operatorname{Spec}(A)$
en los puntos maximales de $\operatorname{Spec}(A)$ son regulares.
\end{proposition}

\begin{proof}
Pongamos $X=\operatorname{Spec}(A)$, $X'=\operatorname{Spec}(\widehat A)$,
y sea $g:X'\to X$ el morfismo canónico. Sea $f:Y\to X$ un morfismo
de desingularización; pongamos $Y'=X'\times_XY$, y sean
$f':Y'\to X'$ y $g':Y'\to Y$ las proyecciones canónicas, de modo que
se tiene un diagrama conmutativo
\[
\xymatrix{
X'\ar[d]_g & Y'\ar[l]_{f'}\ar[d]^{g'}\\
X & Y\ar[l]^f
}
\]

Bastará demostrar que el preesquema $Y'$ es \emph{regular}: en efecto,
como $f$ es birracional y de tipo finito, existe un abierto $U$ denso
en todo $X$ tal que la restricción $f^{-1}(U)\to U$ de $f$ es un
isomorfismo y $f^{-1}(U)$ es denso en todo $Y$ (\textbf{I}, 6.5.5).
Los preesquemas inducidos, respectivamente, en el abierto $g^{-1}(U)$
de $X'$ y en el abierto $g'^{-1}(f^{-1}(U))$ de $Y'$ son, por tanto,
isomorfos; esto demuestra que $g^{-1}(U)$ es regular y, \emph{a fortiori},
también lo son las fibras de $g$ en los puntos maximales de $X$,
que están contenidos en $U$.

\oldpage[IV-2]{219}

Notemos que el morfismo $f'$ es propio (\textbf{II}, 5.4.2) y, en
particular, de tipo finito; por tanto, $Y'$ es noetheriano, puesto que
también lo es $X'$. Sea $a'$ el punto cerrado de $X'$; para ver que $Y'$
es regular, bastará demostrar que, \emph{para todo} $y'\in f'^{-1}(a')$,
\emph{el anillo local $\mathcal O_{Y',y'}$ es regular}. En efecto,
se tendrá entonces $f'^{-1}(a')\subset\operatorname{Reg}(Y')$.
Pero, como $X'$ es el espectro de un anillo local completo y $f'$ es
de tipo finito, $\operatorname{Reg}(Y')$ es \emph{abierto} en $Y'$
(6.12.8); además, como el morfismo $f'$ es \emph{cerrado} y
$f'^{-1}(a')\subset\operatorname{Reg}(Y')$, existe una vecindad $V'$
de $a'$ en $X'$ tal que $f'^{-1}(V')\subset\operatorname{Reg}(Y')$.
Como $\widehat A$ es un anillo local, necesariamente $V'=X'$;
por tanto, $\operatorname{Reg}(Y')=Y'$ y la proposición quedará
demostrada.

Ahora bien, si $a$ es el punto cerrado de $X$, se tiene
$g^{-1}(a)=\{a'\}$ y los cuerpos residuales $\boldsymbol{k}(a)$ y $\boldsymbol{k}(a')$ son
isomorfos; por tanto, $g'^{-1}(f^{-1}(a))=f'^{-1}(a')$, y la restricción
$f'^{-1}(a')\to f^{-1}(a)$ de $g'$ es un \emph{isomorfismo} de estas
dos fibras. Sea $y=g'(y')$ y pongamos $B=\mathcal O_{Y,y}$; si
$\mathfrak n=\mathfrak m_y$ es el ideal maximal de $B$, lo anterior
muestra que en el anillo (noetheriano) $B'=B\otimes_A\widehat A$
existe un único ideal maximal $\mathfrak n'$ sobre $\mathfrak n$ y
se tiene $\mathcal O_{Y',y'}=B'_{\mathfrak n'}$. Para mostrar que el
anillo local $B'_{\mathfrak n'}$ es regular, basta establecer que
también lo es su \emph{completado} (\textbf 0, 17.1.5); ahora bien, por hipótesis
$B$ es regular, y, por tanto, también lo es su completado $\widehat B$.
La conclusión se seguirá, pues, del lema siguiente:

\begin{lemma}[7.9.3.1]
\label{IV.2.s7.9.3.1-fr}
Sean $A$, $B$ dos anillos locales noetherianos,
$\rho:A\to B$ un homomorfismo local tal que el anillo
$B'=B\otimes_A\widehat A$ sea noetheriano. Entonces, para todo ideal
maximal $\mathfrak n'$ de $B'$ sobre el ideal maximal $\mathfrak n$
de $B$ y el ideal maximal de $\widehat A$, los completados de $B$
y de $B'_{\mathfrak n'}$ son isomorfos.
\end{lemma}

\begin{proof}
Sea $\mathfrak m$ el ideal maximal de $A$ y pongamos
$B''=B'_{\mathfrak n'}$. Para todo entero $h>0$ se tiene
$B'/\mathfrak m^hB'=B\otimes_A(\widehat A/\mathfrak m^h\widehat A)$;
pero $\widehat A/\mathfrak m^h\widehat A$ es isomorfo a
$A/\mathfrak m^h$, por lo que $B'/\mathfrak m^hB'$ es isomorfo a
$B/\mathfrak m^hB$ y, en particular, es un anillo local cuyo ideal
maximal es $\mathfrak n'/\mathfrak m^hB'$. Por tanto,
$B''/\mathfrak m^hB''$, isomorfo a
$(B'/\mathfrak m^hB')_{\mathfrak n'}$, es $B$-isomorfo a
$B'/\mathfrak m^hB'$ y, finalmente, a $B/\mathfrak m^hB$.
En particular, el ideal maximal de $B''$ es igual a $\mathfrak nB''$,
y $B''/\mathfrak n^hB''$ es isomorfo a $B/\mathfrak n^h$, por lo
anterior. Como $\widehat{B}''=\varprojlim_h(B''/\mathfrak n^hB'')$,
queda demostrado el lema y, con él, (7.9.3).
\end{proof}
\end{proof}
